% --- Archon physics formalization source begin ---
% archon:physics
% archon:covers PhyXMiniProblems/problem_phyx_mini_0719.lean
% archon:source-report reports/phyx_mini/problem_phyx_mini_0719.source.json

\chapter{Physics problem phyx\_mini\_0719}
\label{ch:PhyXMiniProblems_problem_phyx_mini_0719}

\paragraph{Problem source.}
\#\# Physical scenario

A pendulum consists of a massless, rigid rod with a mass at one end. The other end is pivoted on a frictionless pivot so that the rod can rotate in a complete circle. The pendulum is inverted, with the mass directly above the pivot point, then released. The speed of the mass as it passes through the lowest point is \textbackslash{}(5.0\textbackslash{},\textbackslash{}mathrm\{m/s\}\textbackslash{}). The pendulum later undergoes small-amplitude oscillations at the bottom of the arc.

\#\# Question

What will its frequency be?

\#\# Answer choices

- A: A: 0.42Hz

- B: B: 0.52Hz

- C: C: 0.62Hz

- D: D: 0.72Hz

\#\# Figure caption (auxiliary; use the image as primary evidence)

The image depicts a pendulum-like system with key components and information annotated. 

- A circular path is shown, with a rod or string labeled "L" attached to a pivot point at the top.
- There are two positions marked on the circle: the highest point and the lowest point.
- At the highest point, there's a label indicating a "Before" state with \textbackslash{}( y\_i = 2L \textbackslash{}) and \textbackslash{}( v\_i = 0 \textbackslash{}, \textbackslash{}text\{m/s\} \textbackslash{}), suggesting initial height and velocity.
- At the lowest point, a label indicates an "After" state with \textbackslash{}( y\_f = 0 \textbackslash{}, \textbackslash{}text\{m\} \textbackslash{}) and \textbackslash{}( v\_f = 5.0 \textbackslash{}, \textbackslash{}text\{m/s\} \textbackslash{}), for final height and velocity.
- The pivot is marked clearly with a note, connected to the bottom sphere by the line "L".
- Arrows along the circle indicate the direction of motion from the highest point down to the lowest point.

\#\# Dataset metadata

Rotational Motion; Physical Model Grounding Reasoning, Multi-Formula Reasoning

\paragraph{Recorded answer/context.}
C: C: 0.62Hz

\paragraph{Figure/image path.}
/root/proposal\_for\_physic/science-mango/phyx\_mini\_run/phyx\_data/test\_image/719.png

\paragraph{Formalization target.}
create a compiling Lean file with sorry bodies at `PhyXMiniProblems/problem\_phyx\_mini\_0719.lean`. The Lean declarations must preserve the physical quantities, dimensions or dimensional roles, figure labels, governing-law hypotheses, and final relation expressed by this problem.
Use Mathlib/Physlib names found through LeanExplore where available. If a physics API is missing, introduce faithful local abstractions rather than scalar placeholder aliases.

\begin{theorem}[Physics formalization target]
\label{thm:physics:phyx_mini_0719:target}
\lean{PhyXMiniProblems.ProblemPhyXMini0719.invertedRodPendulumFrequency_is_answer_C}
\uses{def:physics:phyx-mini-0719:phyxminiproblems-problemphyxmini0719-frequencyinhertz, def:physics:phyx-mini-0719:phyxminiproblems-problemphyxmini0719-invertedrodpendulumsetup, def:physics:phyx-mini-0719:phyxminiproblems-problemphyxmini0719-matchesprimaryfigure, def:physics:phyx-mini-0719:phyxminiproblems-problemphyxmini0719-matchesproblemdata, def:physics:phyx-mini-0719:phyxminiproblems-problemphyxmini0719-usesstandardearthgravity, def:physics:phyx-mini-0719:phyxminiproblems-problemphyxmini0719-hasphysicalparameters, def:physics:phyx-mini-0719:phyxminiproblems-problemphyxmini0719-satisfiesinvertedpendulumlaws, def:physics:phyx-mini-0719:phyxminiproblems-problemphyxmini0719-answerfrequencyinhertz, def:physics:phyx-mini-0719:phyxminiproblems-problemphyxmini0719-roundstodisplayedhundredth, def:physics:phyx-mini-0719:phyxminiproblems-problemphyxmini0719-isuniqueclosestanswerchoice}
The assigned autoformalize agent should translate this physics problem into Lean declarations in the covered file, with theorem and lemma proof bodies written as `by sorry`.
\end{theorem}
\begin{proof}
This is an autoformalization task, not a proof task. Produce statements that can later be proved by the physics prover without weakening the physical model.
\end{proof}

% --- Archon declaration topology begin ---
\section{Declaration topology}
\begin{definition}[acceleration Dimension]
  \label{def:physics:phyx-mini-0719:phyxminiproblems-problemphyxmini0719-accelerationdimension}
  \lean{PhyXMiniProblems.ProblemPhyXMini0719.accelerationDimension}
  Acceleration has physical dimension length divided by time squared
\end{definition}
\begin{proof}
  This is the declared model definition of acceleration Dimension.
\end{proof}
\begin{definition}[Mass Quantity]
  \label{def:physics:phyx-mini-0719:phyxminiproblems-problemphyxmini0719-massquantity}
  \lean{PhyXMiniProblems.ProblemPhyXMini0719.MassQuantity}
  A nonnegative physical mass, independent of the unit used to read it
\end{definition}
\begin{proof}
  This is the declared model definition of Mass Quantity.
\end{proof}
\begin{definition}[Length Quantity]
  \label{def:physics:phyx-mini-0719:phyxminiproblems-problemphyxmini0719-lengthquantity}
  \lean{PhyXMiniProblems.ProblemPhyXMini0719.LengthQuantity}
  A nonnegative physical length, independent of the unit used to read it
\end{definition}
\begin{proof}
  This is the declared model definition of Length Quantity.
\end{proof}
\begin{definition}[Speed Quantity]
  \label{def:physics:phyx-mini-0719:phyxminiproblems-problemphyxmini0719-speedquantity}
  \lean{PhyXMiniProblems.ProblemPhyXMini0719.SpeedQuantity}
  A nonnegative physical speed, with dimension length divided by time
\end{definition}
\begin{proof}
  This is the declared model definition of Speed Quantity.
\end{proof}
\begin{definition}[Acceleration Quantity]
  \label{def:physics:phyx-mini-0719:phyxminiproblems-problemphyxmini0719-accelerationquantity}
  \lean{PhyXMiniProblems.ProblemPhyXMini0719.AccelerationQuantity}
  \uses{def:physics:phyx-mini-0719:phyxminiproblems-problemphyxmini0719-accelerationdimension}
  A nonnegative physical acceleration magnitude
\end{definition}
\begin{proof}
  This is the declared model definition of Acceleration Quantity.
\end{proof}
\begin{definition}[Frequency Quantity]
  \label{def:physics:phyx-mini-0719:phyxminiproblems-problemphyxmini0719-frequencyquantity}
  \lean{PhyXMiniProblems.ProblemPhyXMini0719.FrequencyQuantity}
  Cyclic frequency, with physical dimension inverse time
\end{definition}
\begin{proof}
  This is the declared model definition of Frequency Quantity.
\end{proof}
\begin{definition}[Angular Frequency Quantity]
  \label{def:physics:phyx-mini-0719:phyxminiproblems-problemphyxmini0719-angularfrequencyquantity}
  \lean{PhyXMiniProblems.ProblemPhyXMini0719.AngularFrequencyQuantity}
  Angular frequency; radians are dimensionless, so its dimension is inverse time
\end{definition}
\begin{proof}
  This is the declared model definition of Angular Frequency Quantity.
\end{proof}
\begin{definition}[quantity SIReadout]
  \label{def:physics:phyx-mini-0719:phyxminiproblems-problemphyxmini0719-quantitysireadout}
  \lean{PhyXMiniProblems.ProblemPhyXMini0719.quantitySIReadout}
  Read a nonnegative dimensionful quantity in coherent SI units
\end{definition}
\begin{proof}
  This is the declared model definition of quantity SIReadout.
\end{proof}
\begin{definition}[mass In Kilograms]
  \label{def:physics:phyx-mini-0719:phyxminiproblems-problemphyxmini0719-massinkilograms}
  \lean{PhyXMiniProblems.ProblemPhyXMini0719.massInKilograms}
  \uses{def:physics:phyx-mini-0719:phyxminiproblems-problemphyxmini0719-massquantity, def:physics:phyx-mini-0719:phyxminiproblems-problemphyxmini0719-quantitysireadout}
  Kilogram readout of a physical mass
\end{definition}
\begin{proof}
  This is the declared model definition of mass In Kilograms.
\end{proof}
\begin{definition}[length In Meters]
  \label{def:physics:phyx-mini-0719:phyxminiproblems-problemphyxmini0719-lengthinmeters}
  \lean{PhyXMiniProblems.ProblemPhyXMini0719.lengthInMeters}
  \uses{def:physics:phyx-mini-0719:phyxminiproblems-problemphyxmini0719-lengthquantity, def:physics:phyx-mini-0719:phyxminiproblems-problemphyxmini0719-quantitysireadout}
  Meter readout of a physical length or height
\end{definition}
\begin{proof}
  This is the declared model definition of length In Meters.
\end{proof}
\begin{definition}[speed In Meters Per Second]
  \label{def:physics:phyx-mini-0719:phyxminiproblems-problemphyxmini0719-speedinmeterspersecond}
  \lean{PhyXMiniProblems.ProblemPhyXMini0719.speedInMetersPerSecond}
  \uses{def:physics:phyx-mini-0719:phyxminiproblems-problemphyxmini0719-speedquantity, def:physics:phyx-mini-0719:phyxminiproblems-problemphyxmini0719-quantitysireadout}
  Meter-per-second readout of a physical speed
\end{definition}
\begin{proof}
  This is the declared model definition of speed In Meters Per Second.
\end{proof}
\begin{definition}[acceleration In Meters Per Second Squared]
  \label{def:physics:phyx-mini-0719:phyxminiproblems-problemphyxmini0719-accelerationinmeterspersecondsquared}
  \lean{PhyXMiniProblems.ProblemPhyXMini0719.accelerationInMetersPerSecondSquared}
  \uses{def:physics:phyx-mini-0719:phyxminiproblems-problemphyxmini0719-accelerationquantity, def:physics:phyx-mini-0719:phyxminiproblems-problemphyxmini0719-quantitysireadout}
  Meter-per-second-squared readout of an acceleration magnitude
\end{definition}
\begin{proof}
  This is the declared model definition of acceleration In Meters Per Second Squared.
\end{proof}
\begin{definition}[frequency In Hertz]
  \label{def:physics:phyx-mini-0719:phyxminiproblems-problemphyxmini0719-frequencyinhertz}
  \lean{PhyXMiniProblems.ProblemPhyXMini0719.frequencyInHertz}
  \uses{def:physics:phyx-mini-0719:phyxminiproblems-problemphyxmini0719-frequencyquantity, def:physics:phyx-mini-0719:phyxminiproblems-problemphyxmini0719-quantitysireadout}
  Hertz readout of a cyclic frequency
\end{definition}
\begin{proof}
  This is the declared model definition of frequency In Hertz.
\end{proof}
\begin{definition}[angular Frequency In Radians Per Second]
  \label{def:physics:phyx-mini-0719:phyxminiproblems-problemphyxmini0719-angularfrequencyinradianspersecond}
  \lean{PhyXMiniProblems.ProblemPhyXMini0719.angularFrequencyInRadiansPerSecond}
  \uses{def:physics:phyx-mini-0719:phyxminiproblems-problemphyxmini0719-angularfrequencyquantity, def:physics:phyx-mini-0719:phyxminiproblems-problemphyxmini0719-quantitysireadout}
  Radian-per-second readout of an angular frequency
\end{definition}
\begin{proof}
  This is the declared model definition of angular Frequency In Radians Per Second.
\end{proof}
\begin{definition}[Path Position]
  \label{def:physics:phyx-mini-0719:phyxminiproblems-problemphyxmini0719-pathposition}
  \lean{PhyXMiniProblems.ProblemPhyXMini0719.PathPosition}
  Distinguished bob positions on the circular path
\end{definition}
\begin{proof}
  This is the declared model definition of Path Position.
\end{proof}
\begin{definition}[Figure Label]
  \label{def:physics:phyx-mini-0719:phyxminiproblems-problemphyxmini0719-figurelabel}
  \lean{PhyXMiniProblems.ProblemPhyXMini0719.FigureLabel}
  Text and mathematical labels visible in the primary image
\end{definition}
\begin{proof}
  This is the declared model definition of Figure Label.
\end{proof}
\begin{definition}[Figure Anchor]
  \label{def:physics:phyx-mini-0719:phyxminiproblems-problemphyxmini0719-figureanchor}
  \lean{PhyXMiniProblems.ProblemPhyXMini0719.FigureAnchor}
  Anchors to which the handwritten labels refer
\end{definition}
\begin{proof}
  This is the declared model definition of Figure Anchor.
\end{proof}
\begin{definition}[Inverted Pendulum Figure]
  \label{def:physics:phyx-mini-0719:phyxminiproblems-problemphyxmini0719-invertedpendulumfigure}
  \lean{PhyXMiniProblems.ProblemPhyXMini0719.InvertedPendulumFigure}
  \uses{def:physics:phyx-mini-0719:phyxminiproblems-problemphyxmini0719-pathposition, def:physics:phyx-mini-0719:phyxminiproblems-problemphyxmini0719-figurelabel, def:physics:phyx-mini-0719:phyxminiproblems-problemphyxmini0719-figureanchor}
  Qualitative evidence read directly from the supplied bitmap. Numerical readouts are attached to the dimensionful quantities in MatchesPrimaryFigure and MatchesProblemData below
\end{definition}
\begin{proof}
  This is the declared model definition of Inverted Pendulum Figure.
\end{proof}
\begin{definition}[Bob Model]
  \label{def:physics:phyx-mini-0719:phyxminiproblems-problemphyxmini0719-bobmodel}
  \lean{PhyXMiniProblems.ProblemPhyXMini0719.BobModel}
  Idealized model of the body attached to the moving end of the rod
\end{definition}
\begin{proof}
  This is the declared model definition of Bob Model.
\end{proof}
\begin{definition}[Rod Model]
  \label{def:physics:phyx-mini-0719:phyxminiproblems-problemphyxmini0719-rodmodel}
  \lean{PhyXMiniProblems.ProblemPhyXMini0719.RodModel}
  Mechanical model assigned to the pendulum rod
\end{definition}
\begin{proof}
  This is the declared model definition of Rod Model.
\end{proof}
\begin{definition}[Pivot Model]
  \label{def:physics:phyx-mini-0719:phyxminiproblems-problemphyxmini0719-pivotmodel}
  \lean{PhyXMiniProblems.ProblemPhyXMini0719.PivotModel}
  Dissipation model assigned to the central pivot
\end{definition}
\begin{proof}
  This is the declared model definition of Pivot Model.
\end{proof}
\begin{definition}[Path Model]
  \label{def:physics:phyx-mini-0719:phyxminiproblems-problemphyxmini0719-pathmodel}
  \lean{PhyXMiniProblems.ProblemPhyXMini0719.PathModel}
  Extent of the trajectory allowed by the rod and pivot
\end{definition}
\begin{proof}
  This is the declared model definition of Path Model.
\end{proof}
\begin{definition}[Oscillation Regime]
  \label{def:physics:phyx-mini-0719:phyxminiproblems-problemphyxmini0719-oscillationregime}
  \lean{PhyXMiniProblems.ProblemPhyXMini0719.OscillationRegime}
  Approximation regime for the later motion about the bottom of the arc
\end{definition}
\begin{proof}
  This is the declared model definition of Oscillation Regime.
\end{proof}
\begin{definition}[Inverted Rod Pendulum Setup]
  \label{def:physics:phyx-mini-0719:phyxminiproblems-problemphyxmini0719-invertedrodpendulumsetup}
  \lean{PhyXMiniProblems.ProblemPhyXMini0719.InvertedRodPendulumSetup}
  \uses{def:physics:phyx-mini-0719:phyxminiproblems-problemphyxmini0719-massquantity, def:physics:phyx-mini-0719:phyxminiproblems-problemphyxmini0719-lengthquantity, def:physics:phyx-mini-0719:phyxminiproblems-problemphyxmini0719-speedquantity, def:physics:phyx-mini-0719:phyxminiproblems-problemphyxmini0719-accelerationquantity, def:physics:phyx-mini-0719:phyxminiproblems-problemphyxmini0719-frequencyquantity, def:physics:phyx-mini-0719:phyxminiproblems-problemphyxmini0719-angularfrequencyquantity, def:physics:phyx-mini-0719:phyxminiproblems-problemphyxmini0719-pathposition, def:physics:phyx-mini-0719:phyxminiproblems-problemphyxmini0719-invertedpendulumfigure, def:physics:phyx-mini-0719:phyxminiproblems-problemphyxmini0719-bobmodel, def:physics:phyx-mini-0719:phyxminiproblems-problemphyxmini0719-rodmodel, def:physics:phyx-mini-0719:phyxminiproblems-problemphyxmini0719-pivotmodel, def:physics:phyx-mini-0719:phyxminiproblems-problemphyxmini0719-pathmodel, def:physics:phyx-mini-0719:phyxminiproblems-problemphyxmini0719-oscillationregime}
  Independent physical quantities and modeling choices for the experiment. The rod length and both frequencies are unconstrained fields: none is defined to be the requested numerical answer
\end{definition}
\begin{proof}
  This is the declared model definition of Inverted Rod Pendulum Setup.
\end{proof}
\begin{definition}[Matches Primary Figure]
  \label{def:physics:phyx-mini-0719:phyxminiproblems-problemphyxmini0719-matchesprimaryfigure}
  \lean{PhyXMiniProblems.ProblemPhyXMini0719.MatchesPrimaryFigure}
  \uses{def:physics:phyx-mini-0719:phyxminiproblems-problemphyxmini0719-invertedrodpendulumsetup}
  The primary-image readout. It places Before, y\_i, and v\_i at the highest point, After, y\_f, and v\_f at the lowest point, L on the lowest radius, and identifies the central pivot. These are geometric and diagrammatic facts, not frequency conclusions
\end{definition}
\begin{proof}
  This is the declared model definition of Matches Primary Figure.
\end{proof}
\begin{definition}[Matches Problem Data]
  \label{def:physics:phyx-mini-0719:phyxminiproblems-problemphyxmini0719-matchesproblemdata}
  \lean{PhyXMiniProblems.ProblemPhyXMini0719.MatchesProblemData}
  \uses{def:physics:phyx-mini-0719:phyxminiproblems-problemphyxmini0719-lengthinmeters, def:physics:phyx-mini-0719:phyxminiproblems-problemphyxmini0719-speedinmeterspersecond, def:physics:phyx-mini-0719:phyxminiproblems-problemphyxmini0719-invertedrodpendulumsetup}
  Numerical and qualitative data stated in the prose and written on the image. The lowest point is the zero of height, the inverted release point is 2 L above it, and the measured bottom speed is 5.0 m/s. No frequency value or answer choice is asserted here
\end{definition}
\begin{proof}
  This is the declared model definition of Matches Problem Data.
\end{proof}
\begin{definition}[Uses Standard Earth Gravity]
  \label{def:physics:phyx-mini-0719:phyxminiproblems-problemphyxmini0719-usesstandardearthgravity}
  \lean{PhyXMiniProblems.ProblemPhyXMini0719.UsesStandardEarthGravity}
  \uses{def:physics:phyx-mini-0719:phyxminiproblems-problemphyxmini0719-accelerationinmeterspersecondsquared, def:physics:phyx-mini-0719:phyxminiproblems-problemphyxmini0719-invertedrodpendulumsetup}
  The standard near-Earth textbook calibration implicit in the recorded multiple-choice answer. The source does not print a value of g, so this is kept separate from the image and prose readouts
\end{definition}
\begin{proof}
  This is the declared model definition of Uses Standard Earth Gravity.
\end{proof}
\begin{definition}[Has Physical Parameters]
  \label{def:physics:phyx-mini-0719:phyxminiproblems-problemphyxmini0719-hasphysicalparameters}
  \lean{PhyXMiniProblems.ProblemPhyXMini0719.HasPhysicalParameters}
  \uses{def:physics:phyx-mini-0719:phyxminiproblems-problemphyxmini0719-massinkilograms, def:physics:phyx-mini-0719:phyxminiproblems-problemphyxmini0719-lengthinmeters, def:physics:phyx-mini-0719:phyxminiproblems-problemphyxmini0719-accelerationinmeterspersecondsquared, def:physics:phyx-mini-0719:phyxminiproblems-problemphyxmini0719-frequencyinhertz, def:physics:phyx-mini-0719:phyxminiproblems-problemphyxmini0719-angularfrequencyinradianspersecond, def:physics:phyx-mini-0719:phyxminiproblems-problemphyxmini0719-invertedrodpendulumsetup}
  Positivity and nondegeneracy conditions for the physical branch
\end{definition}
\begin{proof}
  This is the declared model definition of Has Physical Parameters.
\end{proof}
\begin{definition}[pendulum Oscillator SI]
  \label{def:physics:phyx-mini-0719:phyxminiproblems-problemphyxmini0719-pendulumoscillatorsi}
  \lean{PhyXMiniProblems.ProblemPhyXMini0719.pendulumOscillatorSI}
  \uses{def:physics:phyx-mini-0719:phyxminiproblems-problemphyxmini0719-massinkilograms, def:physics:phyx-mini-0719:phyxminiproblems-problemphyxmini0719-lengthinmeters, def:physics:phyx-mini-0719:phyxminiproblems-problemphyxmini0719-accelerationinmeterspersecondsquared, def:physics:phyx-mini-0719:phyxminiproblems-problemphyxmini0719-invertedrodpendulumsetup, def:physics:phyx-mini-0719:phyxminiproblems-problemphyxmini0719-hasphysicalparameters}
  The small-angle pendulum as a generalized harmonic oscillator. Its generalized inertia is m L², while the linearized gravitational restoring coefficient is m g L; their ratio is therefore g / L
\end{definition}
\begin{proof}
  This is the declared model definition of pendulum Oscillator SI.
\end{proof}
\begin{definition}[pendulum Linear Coefficient SI]
  \label{def:physics:phyx-mini-0719:phyxminiproblems-problemphyxmini0719-pendulumlinearcoefficientsi}
  \lean{PhyXMiniProblems.ProblemPhyXMini0719.pendulumLinearCoefficientSI}
  \uses{def:physics:phyx-mini-0719:phyxminiproblems-problemphyxmini0719-lengthinmeters, def:physics:phyx-mini-0719:phyxminiproblems-problemphyxmini0719-accelerationinmeterspersecondsquared, def:physics:phyx-mini-0719:phyxminiproblems-problemphyxmini0719-invertedrodpendulumsetup}
  The coefficient g / L in the equation linearized at the bottom equilibrium
\end{definition}
\begin{proof}
  This is the declared model definition of pendulum Linear Coefficient SI.
\end{proof}
\begin{definition}[Has Valid Pendulum Linearization]
  \label{def:physics:phyx-mini-0719:phyxminiproblems-problemphyxmini0719-hasvalidpendulumlinearization}
  \lean{PhyXMiniProblems.ProblemPhyXMini0719.HasValidPendulumLinearization}
  \uses{def:physics:phyx-mini-0719:phyxminiproblems-problemphyxmini0719-invertedrodpendulumsetup, def:physics:phyx-mini-0719:phyxminiproblems-problemphyxmini0719-pendulumlinearcoefficientsi}
  An explicit local/asymptotic contract for the small-angle approximation. The exact restoring acceleration vanishes at the bottom, has derivative -g / L there, and differs from its linear part by o(θ) as θ → 0. Thus no equality for a finite-amplitude nonlinear oscillation is asserted
\end{definition}
\begin{proof}
  This is the declared model definition of Has Valid Pendulum Linearization.
\end{proof}
\begin{definition}[Satisfies Inverted Pendulum Laws]
  \label{def:physics:phyx-mini-0719:phyxminiproblems-problemphyxmini0719-satisfiesinvertedpendulumlaws}
  \lean{PhyXMiniProblems.ProblemPhyXMini0719.SatisfiesInvertedPendulumLaws}
  \uses{def:physics:phyx-mini-0719:phyxminiproblems-problemphyxmini0719-massinkilograms, def:physics:phyx-mini-0719:phyxminiproblems-problemphyxmini0719-lengthinmeters, def:physics:phyx-mini-0719:phyxminiproblems-problemphyxmini0719-speedinmeterspersecond, def:physics:phyx-mini-0719:phyxminiproblems-problemphyxmini0719-accelerationinmeterspersecondsquared, def:physics:phyx-mini-0719:phyxminiproblems-problemphyxmini0719-frequencyinhertz, def:physics:phyx-mini-0719:phyxminiproblems-problemphyxmini0719-angularfrequencyinradianspersecond, def:physics:phyx-mini-0719:phyxminiproblems-problemphyxmini0719-invertedrodpendulumsetup, def:physics:phyx-mini-0719:phyxminiproblems-problemphyxmini0719-hasphysicalparameters, def:physics:phyx-mini-0719:phyxminiproblems-problemphyxmini0719-pendulumoscillatorsi, def:physics:phyx-mini-0719:phyxminiproblems-problemphyxmini0719-pendulumlinearcoefficientsi, def:physics:phyx-mini-0719:phyxminiproblems-problemphyxmini0719-hasvalidpendulumlinearization}
  The governing laws used to infer the rod length and then the linearized frequency prediction. The first field is conservation of m g y + 1/2 m v² between the two pictured states. The second is the exact nonlinear sine restoring law. The final fields use Physlib's harmonic-oscillator angular frequency only for the model obtained by local linearization, and the general conversion ω = 2πf. None substitutes 5.0 m/s, assigns a numerical length, or states the requested 0.62 Hz conclusion
\end{definition}
\begin{proof}
  This is the declared model definition of Satisfies Inverted Pendulum Laws.
\end{proof}
\begin{theorem}[exact Pendulum Law has Valid Local Linearization]
  \label{thm:physics:phyx-mini-0719:phyxminiproblems-problemphyxmini0719-exactpendulumlaw-hasvalidlocallinearization}
  \lean{PhyXMiniProblems.ProblemPhyXMini0719.exactPendulumLaw_hasValidLocalLinearization}
  \uses{def:physics:phyx-mini-0719:phyxminiproblems-problemphyxmini0719-invertedrodpendulumsetup, def:physics:phyx-mini-0719:phyxminiproblems-problemphyxmini0719-hasphysicalparameters, def:physics:phyx-mini-0719:phyxminiproblems-problemphyxmini0719-hasvalidpendulumlinearization, def:physics:phyx-mini-0719:phyxminiproblems-problemphyxmini0719-satisfiesinvertedpendulumlaws}
  The exact sine law supplies the local derivative/remainder contract required to pass from the nonlinear pendulum to its bottom-equilibrium linearization. This is a mathematical consequence of the general governing law, not a numerical frequency premise
\end{theorem}
\begin{proof}
  The conclusion follows from the governing relations and figure readouts named in the statement of exact Pendulum Law has Valid Local Linearization.
\end{proof}
\begin{definition}[Answer Choice]
  \label{def:physics:phyx-mini-0719:phyxminiproblems-problemphyxmini0719-answerchoice}
  \lean{PhyXMiniProblems.ProblemPhyXMini0719.AnswerChoice}
  Labels of the four answer choices in the source
\end{definition}
\begin{proof}
  This is the declared model definition of Answer Choice.
\end{proof}
\begin{definition}[answer Frequency In Hertz]
  \label{def:physics:phyx-mini-0719:phyxminiproblems-problemphyxmini0719-answerfrequencyinhertz}
  \lean{PhyXMiniProblems.ProblemPhyXMini0719.answerFrequencyInHertz}
  \uses{def:physics:phyx-mini-0719:phyxminiproblems-problemphyxmini0719-answerchoice}
  Frequency in hertz printed beside each source answer choice
\end{definition}
\begin{proof}
  This is the declared model definition of answer Frequency In Hertz.
\end{proof}
\begin{definition}[Rounds To Displayed Hundredth]
  \label{def:physics:phyx-mini-0719:phyxminiproblems-problemphyxmini0719-roundstodisplayedhundredth}
  \lean{PhyXMiniProblems.ProblemPhyXMini0719.RoundsToDisplayedHundredth}
  A computed frequency rounds to a displayed hundredth of a hertz
\end{definition}
\begin{proof}
  This is the declared model definition of Rounds To Displayed Hundredth.
\end{proof}
\begin{definition}[Is Unique Closest Answer Choice]
  \label{def:physics:phyx-mini-0719:phyxminiproblems-problemphyxmini0719-isuniqueclosestanswerchoice}
  \lean{PhyXMiniProblems.ProblemPhyXMini0719.IsUniqueClosestAnswerChoice}
  \uses{def:physics:phyx-mini-0719:phyxminiproblems-problemphyxmini0719-answerchoice, def:physics:phyx-mini-0719:phyxminiproblems-problemphyxmini0719-answerfrequencyinhertz}
  A displayed answer is strictly closer than every other listed choice
\end{definition}
\begin{proof}
  This is the declared model definition of Is Unique Closest Answer Choice.
\end{proof}
% --- Archon declaration topology end ---
% --- Archon physics formalization source end ---
